\documentclass[../../main.tex]{subfiles}
\graphicspath{{\subfix{../../image/}}}

\begin{document}

\section{习题}\label{section:ch01-exercises}

\begin{exercise}\label{exercise:ch01-分类}
将 \S1.2 中的各偏微分方程分类如下：
\begin{enumerate}[(a)]
\item 该 PDE 是线性、半线性、拟线性还是完全非线性的？
\item 该 PDE 的阶是多少？
\end{enumerate}
\end{exercise}

\begin{exercise}\label{exercise:ch01-导数个数}
设 $k$ 为正整数。证明：定义在 $\mathbb{R}^n$ 上的光滑函数一般有
\begin{align*}
\binom{n+k-1}{k}=\binom{n+k-1}{n-1}
\end{align*}
个互不相同的 $k$ 阶偏导数。
（提示：这相当于在一行 $k$ 个符号 $\circ$ 中插入 $n-1$ 个分隔符 $|$ 的方式数；例如 $\circ\circ\,||\,\circ\circ\circ\,|\,\circ\,|\,\circ\circ\circ\,||\,\circ\circ\circ\circ\,|$。请解释为何每一种这样的排列恰好对应于一个 $k$ 阶偏导数。）
\end{exercise}

接下来的几道练习用于熟悉\cref{chapter:appA-记号} 中引入的多重指标记号。

\begin{exercise}\label{exercise:ch01-多项式定理}
证明多项式定理：
\begin{align*}
(x_1+\cdots+x_n)^k=\sum_{|\alpha|=k}\binom{|\alpha|}{\alpha}x^\alpha,
\end{align*}
其中 $\binom{|\alpha|}{\alpha}\triangleq\dfrac{|\alpha|!}{\alpha!}$，$\alpha!=\alpha_1!\alpha_2!\cdots\alpha_n!$，$x^\alpha=x_1^{\alpha_1}\cdots x_n^{\alpha_n}$。求和遍历所有满足 $|\alpha|=k$ 的多重指标 $\alpha=(\alpha_1,\cdots,\alpha_n)$。
\end{exercise}

\begin{theorem}\label{theorem:ch01-leibniz}
设 \(u,v:\mathbb{R}^n\to\mathbb{R}\) 光滑 , 则对任意多重指标 \(\alpha\) 有
\begin{align*}
D^\alpha(uv)=\sum\limits_{\beta\leqslant\alpha}\binom{\alpha}{\beta}D^\beta u\,D^{\alpha-\beta}v,
\end{align*}
其中 \(\binom{\alpha}{\beta}\triangleq\frac{\alpha!}{\beta!(\alpha-\beta)!}\) , 而 \(\beta\leqslant\alpha\) 表示 \(\beta_i\leqslant\alpha_i\) , \(i=1,\cdots,n\) .
\end{theorem}
\begin{proof}
对 \(|\alpha|\) 作归纳 . 当 \(|\alpha|=0\) 时 , \(\alpha=0\) , 公式为
\begin{align*}
uv=\binom{0}{0}uv,
\end{align*}
显然成立 .

假设对一切 \(|\alpha|\leqslant m\) 公式成立 . 取 \(|\alpha|=m+1\) . 则存在 \(i\in\{1,\cdots,n\}\) 使得 \(\alpha_i\geqslant1\) . 令
\begin{align*}
\alpha'\triangleq\alpha-e_i,
\end{align*}
其中 \(e_i\) 是第 \(i\) 个标准单位向量 . 则 \(|\alpha'|=m\) . 由归纳假设 ,
\begin{align*}
D^{\alpha'}(uv)=\sum\limits_{\beta\leqslant\alpha'}\binom{\alpha'}{\beta}D^\beta u\,D^{\alpha'-\beta}v.
\end{align*}
于是
\begin{align*}
D^\alpha(uv)
&=D_iD^{\alpha'}(uv)=D_i\sum\limits_{\beta\leqslant\alpha'}\binom{\alpha'}{\beta}D^\beta u\,D^{\alpha'-\beta}v\\
&=\sum\limits_{\beta\leqslant\alpha'}\binom{\alpha'}{\beta}\left(D_iD^\beta u\,D^{\alpha'-\beta}v+D^\beta u\,D_iD^{\alpha'-\beta}v\right)\\
&=\sum\limits_{\beta\leqslant\alpha'}\binom{\alpha'}{\beta}D^{\beta+e_i}u\,D^{\alpha'-\beta}v
+\sum\limits_{\beta\leqslant\alpha'}\binom{\alpha'}{\beta}D^\beta u\,D^{\alpha'-\beta+e_i}v.
\end{align*}
在第一个和式中令 \(\gamma\triangleq\beta+e_i\) , 则 \(\beta=\gamma-e_i\) , 且 \(\gamma\leqslant\alpha\) , \(\gamma_i\geqslant1\) . 在第二个和式中令 \(\gamma\triangleq\beta\) , 则 \(\gamma\leqslant\alpha'\) , 即 \(\gamma_i\leqslant\alpha_i-1\) . 因此
\begin{align*}
D^\alpha(uv)
&=\sum\limits_{\substack{\gamma\leqslant\alpha\\ \gamma_i\geqslant1}}\binom{\alpha'}{\gamma-e_i}D^\gamma u\,D^{\alpha-\gamma}v
+\sum\limits_{\substack{\gamma\leqslant\alpha\\ \gamma_i\leqslant\alpha_i-1}}\binom{\alpha'}{\gamma}D^\gamma u\,D^{\alpha-\gamma}v.
\end{align*}
对固定的 \(\gamma\leqslant\alpha\) , 若 \(\gamma_i\geqslant1\) , 则它在第一个和式中出现 ; 若 \(\gamma_i\leqslant\alpha_i-1\) , 则它在第二个和式中出现 . 当 \(1\leqslant\gamma_i\leqslant\alpha_i-1\) 时 , \(\gamma\) 同时在两个和式中出现 , 其系数为
\begin{align*}
\binom{\alpha'}{\gamma-e_i}+\binom{\alpha'}{\gamma}
&=\frac{(\alpha_i-1)!}{(\gamma_i-1)!(\alpha_i-\gamma_i)!}\prod\limits_{j\neq i}\frac{\alpha_j!}{\gamma_j!(\alpha_j-\gamma_j)!}
+\frac{(\alpha_i-1)!}{\gamma_i!(\alpha_i-1-\gamma_i)!}\prod\limits_{j\neq i}\frac{\alpha_j!}{\gamma_j!(\alpha_j-\gamma_j)!}=\binom{\alpha}{\gamma}.
\end{align*}
当 \(\gamma_i=0\) 时 , 只有第二个和式出现 , 系数为
\begin{align*}
\binom{\alpha'}{\gamma}=\binom{\alpha}{\gamma}.
\end{align*}
当 \(\gamma_i=\alpha_i\) 时 , 只有第一个和式出现 , 系数为
\begin{align*}
\binom{\alpha'}{\gamma-e_i}=\binom{\alpha}{\gamma}.
\end{align*}
故合并后得
\begin{align*}
D^\alpha(uv)=\sum\limits_{\beta\leqslant\alpha}\binom{\alpha}{\beta}D^\beta u\,D^{\alpha-\beta}v.
\end{align*}
由归纳原理 , 公式对一切多重指标成立 .
\end{proof}

\begin{theorem}\label{theorem:ch01-taylor}
设 \(f:\mathbb{R}^n\to\mathbb{R}\) 光滑 . 则对每个 \(k=1,2,\cdots\) , 当 \(x\to 0\) 时
\begin{align*}
f(x)=\sum\limits_{|\alpha|\leqslant k}\frac{1}{\alpha!}D^\alpha f(0)x^\alpha+O(|x|^{k+1}).
\end{align*}
\end{theorem}
\begin{proof}
固定 \(x\in\mathbb{R}^n\) , 令
\begin{align*}
g(t)\triangleq f(tx),\quad t\in\mathbb{R}.
\end{align*}
则 \(g\) 光滑 . 由链式法则 , 对 \(j=0,1,\cdots,k+1\) 有
\begin{align*}
g^{(j)}(t)=\sum\limits_{|\alpha|=j}\frac{j!}{\alpha!}D^\alpha f(tx)x^\alpha.
\end{align*}
事实上 , 对 \(j\) 作归纳即可证明 . 当 \(j=0\) 时显然 . 若 \(j\) 时成立 , 则
\begin{align*}
g^{(j+1)}(t)
&=\frac{\mathrm{d}}{\mathrm{d}t}\sum\limits_{|\alpha|=j}\frac{j!}{\alpha!}D^\alpha f(tx)x^\alpha=\sum\limits_{|\alpha|=j}\frac{j!}{\alpha!}\sum\limits_{i=1}^nD_iD^\alpha f(tx)x_ix^\alpha\\
&=\sum\limits_{|\alpha|=j}\sum\limits_{i=1}^n\frac{j!}{\alpha!}D^{\alpha+e_i}f(tx)x^{\alpha+e_i}=\sum\limits_{|\beta|=j+1}\frac{(j+1)!}{\beta!}D^\beta f(tx)x^\beta.
\end{align*}
因此公式成立 .

由一元 Taylor 公式 , 存在积分余项表示
\begin{align*}
g(1)=\sum\limits_{j=0}^k\frac{g^{(j)}(0)}{j!}+\frac{1}{k!}\int_0^1(1-t)^kg^{(k+1)}(t)\mathrm{d}t.
\end{align*}
代入 \(g^{(j)}(0)\) 的表达式得
\begin{align*}
f(x)
&=\sum\limits_{j=0}^k\sum\limits_{|\alpha|=j}\frac{1}{\alpha!}D^\alpha f(0)x^\alpha
+\frac{1}{k!}\int_0^1(1-t)^k\sum\limits_{|\alpha|=k+1}\frac{(k+1)!}{\alpha!}D^\alpha f(tx)x^\alpha\mathrm{d}t\\
&=\sum\limits_{|\alpha|\leqslant k}\frac{1}{\alpha!}D^\alpha f(0)x^\alpha+R_k(x),
\end{align*}
其中
\begin{align*}
R_k(x)\triangleq\frac{(k+1)!}{k!}\sum\limits_{|\alpha|=k+1}\frac{1}{\alpha!}\int_0^1(1-t)^kD^\alpha f(tx)x^\alpha\mathrm{d}t.
\end{align*}
下面估计 \(R_k(x)\) . 因为 \(f\) 光滑 , 所以每个 \(D^\alpha f\) 连续 . 取 \(r>0\) , 当 \(|x|\leqslant r\) 且 \(t\in[0,1]\) 时 , \(tx\) 属于紧集 \(\overline{B(0,r)}\) . 因此存在常数 \(M>0\) 使得
\begin{align*}
|D^\alpha f(tx)|\leqslant M,\quad |\alpha|=k+1,\ t\in[0,1],\ |x|\leqslant r.
\end{align*}
于是
\begin{align*}
|R_k(x)|
&\leqslant\frac{(k+1)!}{k!}\sum\limits_{|\alpha|=k+1}\frac{1}{\alpha!}\int_0^1(1-t)^k|D^\alpha f(tx)|\,|x^\alpha|\mathrm{d}t\\
&\leqslant M\frac{(k+1)!}{k!}\sum\limits_{|\alpha|=k+1}\frac{1}{\alpha!}|x^\alpha|\int_0^1(1-t)^k\mathrm{d}t\\
&\leqslant C\sum\limits_{|\alpha|=k+1}|x^\alpha|,
\end{align*}
其中 \(C\triangleq M\frac{(k+1)!}{k!}\max\limits_{|\alpha|=k+1}\frac{1}{\alpha!}\int_0^1(1-t)^k\mathrm{d}t\) . 又因为对任意多重指标 \(\alpha\) 有
\begin{align*}
|x^\alpha|=|x_1|^{\alpha_1}\cdots|x_n|^{\alpha_n}\leqslant |x|^{|\alpha|},
\end{align*}
所以当 \(|\alpha|=k+1\) 时 ,
\begin{align*}
|x^\alpha|\leqslant |x|^{k+1}.
\end{align*}
因此
\begin{align*}
|R_k(x)|\leqslant C'|x|^{k+1},
\end{align*}
其中 \(C'\) 只依赖于 \(k,n\) 和 \(f\) 在 \(\overline{B(0,r)}\) 上的导数上界 . 故
\begin{align*}
R_k(x)=O(|x|^{k+1}),\quad x\to0.
\end{align*}
定理得证 .
\end{proof}

\end{document}